% 03_quantum_and_signal_dynamics.tex
\chapter{Quantum Formalism and Signal Dynamics}
\label{ch:quantum_signal_dynamics}

\begin{objectivebox}
\begin{itemize}
    \item Map the electromagnetic Lagrangian density to continuous mechanical stress in the vacuum lattice.
    \item Derive the Generalized Uncertainty Principle (GUP) and the Schr\"odinger Equation from discrete signal bandwidth and LC circuit resonance.
    \item Understand wave-particle duality and the physical origin of Quantum Foam as baseline electrical noise.
    \item Derive the physical mechanism of Quantum Entanglement as a topologically protected phase-locked thread on the $\mathcal{M}_A$ lattice, and prove the CHSH violation $|S| = 2\sqrt{2}$ from first principles.
\end{itemize}
\end{objectivebox}

\begin{warningbox}
\textbf{Nomenclature Disclaimer (VCA Bridge):} Throughout the foundational volumes, Standard Model kinematic terms such as \textbf{``Kinetic Energy''} and \textbf{``Potential Energy''} are heavily used as pedagogical stepping stones to bridge academic understanding. However, within the active solvers and true Universal Geometry of the AVE engine, these classical particle terms are purely structural misnomers. 
They explicitly map to the \textbf{Inductive Phase Propagation} (inertia) and the \textbf{Capacitive / Strain Bias} (restoring force) of the $Z_0$ vacuum LC medium. All solver engine logic uses the strict topological limits (\texttt{inductive\_energy} and \texttt{capacitive\_bias}) seamlessly devoid of particle physics analogues.
\end{warningbox}

Standard Quantum Field Theory (QFT) relies on a Lagrangian density ($\mathcal{L}$) describing fields as mathematical operators. In Applied Vacuum Engineering, the continuous quantum formalism is derived from the discrete finite-element signal dynamics of the $\mathcal{M}_A$ lattice.

\section{The Dielectric Lagrangian: Hardware Mechanics}
The substitution of $\xi_{topo}$ converts the electromagnetic Lagrangian density into continuous mechanical stress ($\text{N/m}^2$), grounding Axiom 3 in bulk continuum mechanics. The total macroscopic energy density of the manifold is the sum of the energy stored in the capacitive edges (dielectric strain) and the inductive nodes (kinematic inertia). To construct a relativistically invariant action principle, the Lagrangian difference ($\mathcal{L} = \mathcal{T} - \mathcal{U}$) is evaluated.

The canonical field variable for evaluating transverse waves across a discrete graph is the \textbf{Magnetic Vector Potential} ($\mathbf{A}$), defining the magnetic flux linkage per unit length ($[\text{Wb/m}] = [\text{V}\cdot\text{s/m}]$). Because the generalized velocity of this coordinate is identically the electric field ($\mathbf{E} = -\partial_t \mathbf{A}$), the capacitive energy takes the role of kinetic energy ($\mathcal{T}$), and the inductive energy acts as potential energy ($\mathcal{U}$).
\begin{resultbox}{Dielectric Lagrangian Density}
\begin{equation}
    \mathcal{L}_{AVE} = \frac{1}{2} \epsilon_0 \left| \frac{\partial \mathbf{A}}{\partial t} \right|^2 - \frac{1}{2\mu_0} |\nabla \times \mathbf{A}|^2
\end{equation}
\end{resultbox}

\subsection{Dimensional Proof: The Vector Potential as Mass Flow}
Evaluating the SI dimensions of this continuous field confirms its mechanical identity. Applying the topological conversion constant ($\xi_{topo} \equiv e/\ell_{node}$ measured in $[\text{C/m}]$) to the canonical variable $\mathbf{A}$:
\begin{resultbox}{Vector Potential Dimensions}
\begin{equation}
    [\mathbf{A}] = \left[ \frac{\text{V} \cdot \text{s}}{\text{m}} \right] = \left[ \frac{\text{J} \cdot \text{s}}{\text{C} \cdot \text{m}} \right] = \left[ \frac{\text{kg} \cdot \text{m}^2 \cdot \text{s}}{\text{s}^2 \cdot \text{C} \cdot \text{m}} \right] = \left[ \frac{\text{kg} \cdot \text{m}}{\text{s} \cdot \text{C}} \right]
\end{equation}
\end{resultbox}
By substituting the mathematically exact topological conversion $\text{C} \equiv \xi_{topo} \text{ m}$ derived in Chapter 2, the spatial metric evaluates to:
\begin{resultbox}{Vector Potential as Mass Flow}
\begin{equation}
    [\mathbf{A}] = \left[ \frac{\text{kg} \cdot \text{m}}{\text{s} \cdot (\xi_{topo} \text{ m})} \right] = \mathbf{\xi_{topo}^{-1} \left[ \frac{\text{kg}}{\text{s}} \right]}
\end{equation}
\end{resultbox}
This establishes a fundamental dimensional equivalence: the magnetic vector potential ($\mathbf{A}$) is physically isomorphic to the continuous \textbf{Mass Flow Rate} (linear momentum density) of the vacuum lattice, scaled inversely by the topological dislocation constant. The time derivative therefore carries one extra factor of inverse time: $[\partial_t \mathbf{A}] = \xi_{topo}^{-1}\,[\text{kg/s}^2]$.

When evaluating the full kinetic energy density term using this mechanical substitution,
and retrieving the exact capacitive compliance derivation from Chapter 2 ($\epsilon_{0}\equiv\xi_{topo}^2[\text{N}^{-1}]$),
the topological scaling constants cancel:

\begin{resultbox}{Kinetic Energy Density Dimensions}
\begin{align}
[\mathcal{L}_{kin}] &= \frac{1}{2}\epsilon_{0}|\partial_{t}A|^{2} \nonumber \\
&\Rightarrow (\xi_{topo}^{2}[\text{N}^{-1}])\left(\xi_{topo}^{-1}\frac{\text{kg}}{\text{s}^2}\right)^{2} = \left(\frac{\xi_{topo}^{2}}{\xi_{topo}^{2}}\right)\frac{\text{kg}^{2}}{\text{N}\cdot \text{s}^{4}} \nonumber \\
&= \frac{\text{kg}^{2}}{(\text{kg}\cdot \text{m}/\text{s}^{2})\cdot \text{s}^{4}} \equiv \left[\frac{\text{N}}{\text{m}^{2}}\right]
\end{align}
\end{resultbox}

Minimizing the quantum action is mathematically equivalent to minimizing the continuous
inductive bulk stress (Pascals) of the $\mathcal{M}_{A}$ manifold.

\section{Deriving the Quantum Formalism from Signal Bandwidth}
Standard Quantum Mechanics introduces its formalism---complex Hilbert spaces and non-commuting operators---as axiomatic postulates\cite{feynman1964}. In the AVE framework, these follow as algebraic consequences of transmitting finite-bandwidth signals across a discrete mechanical graph.

\subsection{The Paley-Wiener Hilbert Space}
Because the $\mathcal{M}_A$ lattice has a fundamental pitch $\ell_{node}$, it acts as a spatial Nyquist sampling grid\cite{nyquist1928}. The maximum spatial frequency the lattice can support without aliasing is the Brillouin boundary: $k_{max} = \pi / \ell_{node}$.

By the \textbf{Whittaker-Shannon Interpolation Theorem}, any band-limited signal $\mathbf{A}(\mathbf{x})$ propagating through this discrete lattice can be reconstructed uniquely using a superposition of orthogonal sinc functions\cite{shannon1949}. The set of all such band-limited functions constitutes a Reproducing Kernel Hilbert Space known as the \textbf{Paley-Wiener Space} ($PW_{\pi/\ell_{node}}$).

To map the real-valued physical lattice potential $\mathbf{A}(\mathbf{x},t)$ to the complex continuous quantum state vector $\Psi(\mathbf{x},t)$, the standard signal-processing \textbf{Analytic Signal} representation utilizing the Hilbert Transform ($\mathcal{H}_{transform}$) is applied:
\begin{resultbox}{Analytic Signal Extension}
\begin{equation}
    \Psi(\mathbf{x},t) = \mathbf{A}(\mathbf{x},t) + i \mathcal{H}_{transform}[\mathbf{A}(\mathbf{x},t)]
\end{equation}
\end{resultbox}
The complex continuous Hilbert space of standard quantum mechanics is formally identical to the Paley-Wiener signal-processing representation of the discrete vacuum hardware.

\subsection{The Generalized Uncertainty Principle (GUP)}
On a discrete graph with pitch $\ell_{node}$, continuous coordinate translation is physically impossible. The continuous canonical momentum $\hat{p}_c$ is strictly bounded by the Brillouin zone of the lattice:
\begin{equation}
    p_c \in \left[-\frac{\pi\hbar}{\ell_{node}},\; \frac{\pi\hbar}{\ell_{node}}\right]
\end{equation}
For a macroscopic wave propagating through a stochastic 3D amorphous solid, the effective continuous momentum operator $\langle \hat{P} \rangle$ is defined as an isotropic ensemble average of the symmetric central finite-difference operator across adjacent nodes:
\begin{resultbox}{Continuous Momentum Expectation}
\begin{equation}
    \langle \hat{P} \rangle \approx \frac{\hbar}{\ell_{node}} \sin\left(\frac{\ell_{node} \hat{p}_c}{\hbar}\right)
\end{equation}
\end{resultbox}

Evaluating the exact commutator of the continuous position operator with this discrete lattice momentum ($[\hat{x}, f(\hat{p}_c)] = i\hbar f'(\hat{p}_c)$) yields:
\begin{resultbox}{Discrete Graph Commutator}
\begin{equation}
    [\hat{x}, \langle \hat{P} \rangle] = i\hbar \cos\left(\frac{\ell_{node} \hat{p}_c}{\hbar}\right)
\end{equation}
\end{resultbox}

Applying the generalized Robertson-Schr\"odinger relation yields the rigorous \textbf{Generalized Uncertainty Principle (GUP)} for the discrete vacuum. Because continuous momentum $\Delta x_{SM}$ and the fundamental node spacing are orthogonal hardware constraints, they add in quadrature (Root-Sum-Square):
\begin{resultbox}{The Generalized Uncertainty Principle (GUP)}
\begin{equation}
    \Delta x_{AVE} = \sqrt{(\Delta x_{SM})^2 + \left(\frac{\ell_{node}}{2}\right)^2} \ge \frac{\ell_{node}}{2}
\end{equation}
\end{resultbox}

\textbf{The Physical Origin of the GUP Gap:} In the low-energy limit ($p_c \ll \hbar/\ell_{node}$), the cosine evaluates to $1$, recovering standard Heisenberg physics ($\Delta x \Delta p \ge \hbar/2$). Standard continuum physics assumes that as kinetic momentum increases without bound, $\Delta x$ can shrink arbitrarily. This assumption is the origin of Ultraviolet (UV) divergences in standard Field Theories.

In the AVE framework, as kinetic energies approach the Brillouin zone boundary of the lattice, the cosine expectation value decreases toward zero. The curve separates from the continuum limit and saturates at a finite plateau. This \textbf{GUP Gap} indicates that a pressure wave cannot be localised below the structural pitch of the lattice, providing a built-in UV regularisation.

\begin{examplebox}[Calculating the GUP Cutoff for High-Energy Probes]
\textbf{Problem:} Under the Standard Model, decreasing the position uncertainty $\Delta x$ implies launching particles with increasingly extreme kinetic momentum $\Delta p$. If an attempt is made to resolve the vacuum grid itself, what limits the resolution?

\textbf{Solution:} In standard continuous quantum mechanics, $\Delta x \ge \hbar / (2\Delta p)$. As $\Delta p \to \infty$, theoretical $\Delta x \to 0$. However, in the AVE discrete lattice, the minimum physical wavelength is Nyquist-bounded by the node pitch.
Evaluating the generalized uncertainty formula as continuous $\Delta x_{SM} \to 0$:
\[ \Delta x_{AVE} = \sqrt{(\Delta x_{SM})^2 + \left(\frac{\ell_{node}}{2}\right)^2} \ge \frac{\ell_{node}}{2} \approx 1.93 \times 10^{-13}\,\text{m} \]
Even with infinite continuous momentum, a physical pressure wave cannot resolve structural details smaller than the spatial Nyquist limit of the underlying LC network.
\end{examplebox}

\begin{figure}[h]
    \centering
    \includegraphics[width=1.0\textwidth]{ave_gup_resolution.png}
    \caption{\textbf{The Generalized Uncertainty Principle.} In the continuum limit (red), the uncertainty variance approaches zero, implying unbounded localisation at high energies. In the discrete AVE limit (cyan), the Brillouin boundary forces the finite-difference momentum to plateau, enforcing a minimum localisation length.}
    \label{fig:gup_resolution}
\end{figure}

\subsection{Deriving the Schr\"odinger Equation from Circuit Resonance}
When a topological defect (mass) is synthesized within the graph, it acts as a localized inductive load, imposing a fundamental circuit resonance frequency ($\omega_m = mc^2/\hbar$). This mathematically transforms the massless wave equation into the massive \textbf{Klein-Gordon Equation}:
\begin{resultbox}{Klein-Gordon Equation as Circuit Resonance}
\begin{equation}
    \nabla^2 \mathbf{A} - \frac{1}{c^2}\frac{\partial^2 \mathbf{A}}{\partial t^2} = \left(\frac{mc}{\hbar}\right)^2 \mathbf{A}
\end{equation}
\end{resultbox}

To map this relativistic classical evolution to non-relativistic quantum states, the \textbf{Paraxial Approximation} is applied, factoring out the rest-mass Compton frequency via a slow-varying envelope function $\mathbf{A}(\mathbf{x},t) = \Psi(\mathbf{x},t) e^{-i \omega_m t}$. 

For non-relativistic speeds ($v \ll c$), the second time derivative of the envelope ($\partial_t^2 \Psi$) is negligible. The strict mass resonance terms precisely cancel out:
\begin{resultbox}{The Schr\"odinger Equation}
\begin{equation}
    \nabla^2 \Psi + \frac{2im}{\hbar} \frac{\partial \Psi}{\partial t} = 0 \quad \implies \quad i\hbar \frac{\partial \Psi}{\partial t} = -\frac{\hbar^2}{2m} \nabla^2 \Psi
\end{equation}
\end{resultbox}
The Schr\"odinger Equation evaluates precisely as the paraxial envelope equation of a classical macroscopic pressure wave propagating through the discrete massive $LC$ circuits of the vacuum.

\section{Wave-Particle Duality and The Zero-Impedance Boundary}
The framework asserts that subatomic particles are topological knots where the spatial LC metric reaches dielectric saturation. Two distinct thresholds govern this process:
\begin{itemize}
    \item \textbf{The Kinetic Yield Limit} ($E_k = \sqrt{\alpha} \cdot m_e c^2 \approx 43.65 \text{ keV}$): the threshold where non-linear saturation effects become dominant and the soliton begins to self-confine.
    \item \textbf{The Absolute Snap Limit} ($V_{snap} = m_e c^2 / e \approx 511 \text{ kV}$): the maximum inter-node potential difference before the lattice structurally ruptures (derived in Chapter \ref{ch:macroscopic_moduli}).
\end{itemize}
The kinetic yield governs mass assembly and confinement; the snap limit governs dielectric breakdown and pair production. The Transmission Line mathematics of the saturation boundary provide a derivation of the physical origin of solid matter and wave-particle duality.

\begin{figure}[h!]
    \centering
    \includegraphics[width=1.0\textwidth]{double_slit_sm_vs_ave.png}
    \caption{\textbf{Wave-Particle Duality: SM vs AVE Interpretations.} (Left) The Standard Model models the particle abstractly as a probability wave passing through both slits simultaneously. (Right) The AVE framework models the particle as an objective localized topological defect passing through one slit, while its physical macroscopic hydrodynamic strain (the "Dark Wake") passes through both, causing classic continuous mechanical interference at the detector.}
    \label{fig:double_slit_comparison}
\end{figure}

\subsection{The $0 \ \Omega$ Boundary Condition}
The surrounding relaxed vacuum has a characteristic impedance $Z_{vac} \approx 377 \ \Omega$. At the saturated core of the localised knot, the LC nodes can no longer support alternating transverse displacement, so the effective dynamic RF impedance drops to $0 \ \Omega$ (an RF short circuit).

In transmission line theory, when a wave hits an impedance boundary, the ratio of reflected energy is governed strictly by the Reflection Coefficient ($\Gamma$):
\begin{resultbox}{Transmission Line Reflection}
\begin{equation}
    \Gamma = \frac{Z_{knot} - Z_{vacuum}}{Z_{knot} + Z_{vacuum}}
\end{equation}
\end{resultbox}

By evaluating the ratio at the saturated knot boundary ($Z_{knot} = 0 \ \Omega$):
\begin{resultbox}{The Absolute Impedance Boundary}
\begin{equation}
    \Gamma = \frac{0 - 377}{0 + 377} = \mathbf{-1}
\end{equation}
\end{resultbox}
A Reflection Coefficient of $-1$ signifies total reflectance.

\subsection{Internal Confinement and Matter Assembly}
Because the boundary of the saturated knot represents a steep impedance gradient dropping to $0 \ \Omega$, any acoustic energy circulating \textit{inside} the knot is fully reflected inward. 
A subatomic particle in this model is a stable, self-sustaining acoustic standing wave trapped inside a spherical $0 \ \Omega$ impedance boundary of its own geometric creation.

\subsection{Scattering and The Pauli Exclusion Principle}
Conversely, when an external wave (such as a photon) travels through the $377 \ \Omega$ relaxed vacuum and encounters the knot, it meets the same $0 \ \Omega$ boundary. The photon cannot penetrate the saturated volume; it experiences total reflection. 

This transmission line mismatch provides a mechanical basis for the \textbf{Pauli Exclusion Principle} and the concept of cross-sectional area in particle physics. Two saturated knots cannot occupy the same spatial coordinates because their respective $0 \ \Omega$ boundaries reflect each other's inductive phase energy. Solid matter emerges from vacuum wave mechanics through macroscopic impedance reflection.

\section{The Physical Origin of Quantum Foam and Virtual Particles}
In the standard model of cosmology, the vacuum at the Planck scale is described as a turbulent geometry (``Quantum Foam'') populated by virtual particles that appear and disappear continuously. Standard Quantum Field Theory uses these virtual contributions to balance perturbative equations, which leads to large discrepancies such as the Cosmological Constant Problem, where the predicted vacuum energy density exceeds observations by approximately 120 orders of magnitude.

The AVE framework addresses this discrepancy by replacing virtual particle constructs with the physical dynamics of an active electrical network.

\subsection{Quantum Foam as Baseline RMS Thermal Noise}
Because the physical vacuum $\mathcal{M}_A$ is a literal LC Resonant Network, it is subject to the absolute laws of electrical engineering. In any physical inductor-capacitor (LC) network operating above absolute zero, there exists an irreducible, baseline RMS thermal noise floor. 

What standard physics identifies as ``Quantum Foam''---the underlying geometric turbulence of empty space---is explicitly defined in the AVE framework as the continuous, irreducible electromagnetic AC transients (voltage and current ripples) propagating randomly across the discrete topological grid. It is not geometry itself boiling; it is the chaotic, baseline electrical noise floor of the universe's hardware substrate. This provides a deterministic, continuous mechanical origin for Zero-Point Energy (ZPE) bounded strictly by the finite geometry of the local spatial node.

\subsection{Virtual Particles as Failed Topologies}
In AVE, stable ``Matter'' (such as the electron) is defined as a closed, localised topological knot (specifically, a $0_1$ Unknot---a single closed flux tube loop at minimum ropelength $= 2\pi$) that locks geometrically into the lattice. Maintaining this structural lock requires sustained threshold energy (the $43.65$ keV structural yield limit derived in Volume III, Chapter 1).

When the continuous AC transients (the Quantum Foam) spike, they momentarily twist the local LC phase, creating transient geometric loops. However, because these random spikes generally lack the sustained inductive tension required to close a stable unknot, the intrinsic $\mu_0, \epsilon_0$ tension of the lattice returns the twisted loop to its flat ground state.

In this interpretation, ``Virtual Particles'' are \textbf{failed topologies}: transient, localised phase twists generated by thermal node noise that fail to achieve stable resonant closure, unwinding and dissipating back into the baseline noise floor.

\section{Deterministic Interference and The Measurement Effect}
In the Double Slit Experiment, the topological defect (particle) passes through Slit A, but the continuous transverse inductive wake generated by its motion passes through \textit{both} slits. The particle deterministically navigates the resulting transverse ponderomotive gradients ($\mathbf{F} \propto \nabla |\Psi|^2$) into the quantized standing-wave troughs.

\begin{figure}[h]
    \centering
    \includegraphics[width=1.0\textwidth]{double_slit_decoherence.png}
    \caption{\textbf{Thermodynamic Wavefunction Collapse (Simulation Output).} A continuous PDE finite-difference solver modeling the Double-Slit experiment. \textbf{Case A (Unmeasured):} A classical topological defect translates decisively through Slit 1, while its transverse acoustic wake passes through both slits, generating macroscopic interference fringes ($\nabla |\Psi_{mech}|^2$) across the rear barrier. \textbf{Case B (Measured):} The act of "Measurement" physically requires the insertion of a macroscopic Ohmic detector (a geometric acoustic damper) at Slit 2. The detector physically extracts and thermalizes the acoustic phase energy (Joule friction), deterministically destroying the wave's phase coherence. The quantum "Wavefunction Collapse" cleanly manifests as purely classical structural decoherence, leaving only a single-slit distribution smear.}
    \label{fig:double_slit_wake}
\end{figure}

\subsection{Ohmic Decoherence and the Born Rule}
To measure a quantum state, a macroscopic detector must physically couple to the vacuum lattice. By Axiom 1, any device that couples to the $\mathbf{A}$-field and extracts kinetic energy acts as a resistive mechanical load (where $1 \, \Omega \equiv \xi_{topo}^{-2} \text{ kg/s}$). The physical work extracted into the detector over a measurement interval $\Delta t$ is governed by classical continuous Joule heating ($P = V^2 / R$):
\begin{resultbox}{Extracted Measurement Work}
\begin{equation}
    W_{extracted} = \int P_{load} dt \propto \frac{|\partial_t \mathbf{A}(x_n)|^2}{Z_{detector}} \Delta t
\end{equation}
\end{resultbox}

In a stochastic thermal substrate, the probability that the extracted work triggers a macroscopic discrete event scales identically with the squared amplitude of the local wave envelope.
\begin{resultbox}{Deterministic Born Rule}
\begin{equation}
    P(click | x_n) = \frac{|\partial_t \mathbf{A}(x_n)|^2}{\int |\partial_t \mathbf{A}(\mathbf{x})|^2 d^3x} \equiv |\Psi|^2
\end{equation}
\end{resultbox}
\textbf{The Born Rule} represents the deterministic thermodynamic equation for momentum extraction from a wave-bearing lattice by a thresholded Ohmic load. Placing a detector at Slit B irreversibly thermalizes the spatial pressure wave (decoherence), permanently attenuating the interference gradients.

\section{Non-Linear Dynamics and Topological Shockwaves}
The linear wave equation assumes constant compliance ($\epsilon_0$). Axiom 4, however, defines the vacuum as a non-linear dielectric bounded by the fine-structure limit ($\alpha$). The saturation operator takes a squared geometric form ($n=2$), consistent with QED energy bounds and classical electrodynamics.

To preserve dimensional homogeneity on a 1D continuous transmission line, the telegrapher equations utilize the continuous macroscopic non-linear modulus $\epsilon(\Delta\phi)$:
\begin{resultbox}{Non-Linear Telegrapher Equation}
\begin{equation}
\label{eq:nonlinear_wave}
    \frac{\partial^{2}\Delta\phi}{\partial z^{2}} = \mu_0 \epsilon(\Delta\phi)\frac{\partial^{2}\Delta\phi}{\partial t^{2}} + \mu_0 \frac{d\epsilon}{d\Delta\phi}\left(\frac{\partial \Delta\phi}{\partial t}\right)^{2}
\end{equation}
\end{resultbox}

Enforcing the Saturation Operator defined in Axiom 4:
\begin{resultbox}{Dielectric Saturation Taylor Expansion}
\begin{equation}
    \epsilon(\Delta\phi) = \epsilon_{0}\sqrt{1 - \left(\frac{\Delta\phi}{\alpha}\right)^2} \implies \epsilon(\Delta\phi) \approx \epsilon_0 \left[1 - \frac{1}{2}\left(\frac{\Delta\phi}{\alpha}\right)^2\right]
\end{equation}
\end{resultbox}

The continuous dielectric displacement $D = \epsilon(\Delta\phi) \cdot \Delta\phi$ evaluates to $D_{NL} \approx \epsilon_0 \Delta\phi - \frac{\epsilon_0}{2\alpha^2}(\Delta\phi)^3$. The stored volumetric energy density ($U$) is the integral of the field with respect to displacement ($U = \int \Delta\phi \, dD$):
\begin{resultbox}{Euler-Heisenberg $E^4$ Correction}
\begin{equation}
    U \approx \frac{1}{2}\epsilon_0 (\Delta\phi)^2 - \frac{3}{8\alpha^2}\epsilon_0 (\Delta\phi)^4
\end{equation}
\end{resultbox}

The $(\Delta\phi)^4$ correction term corresponds to the energy density structure of the \textbf{Euler-Heisenberg QED Lagrangian}. The corresponding $D \propto (\Delta\phi)^3$ displacement yields the 3rd-order optical non-linearity associated with the \textbf{Kerr Effect ($\chi^{(3)}$)}.

\begin{figure}[h]
    \centering
    \includegraphics[width=0.95\textwidth]{vacuum_dielectric_saturation.png}
    \caption{\textbf{Axiom 4 Saturation Observables.} (Top) Constitutive permittivity $\varepsilon_{eff} = \varepsilon_0 S(A)$ collapses toward zero as field strain approaches the yield limit. (Bottom) Measurable capacitance $C_{eff} = C_0 / S(A)$ diverges as the same kernel drives the reciprocal response. Three regimes are shaded: \textbf{I} (green) Linear, \textbf{II} (orange) Euler-Heisenberg $E^4$ correction, \textbf{III} (red) full saturation at the Schwinger pair-production threshold.}
    \label{fig:vacuum_dielectric_saturation}
\end{figure}
As the local strain approaches the yield limit, the localised wave speed $c_{eff}(\Delta\phi) = c_0 [1 - (\Delta\phi/\alpha)^2]^{-1/4}$ diverges toward infinity (due to the vanishing permittivity, $\epsilon \to 0$). Because the high-amplitude peak propagates faster than the low-amplitude base, the peak overruns the leading edge, steepening it until it topologically snaps. This forward structural shockwave provides a continuous mechanical origin for discrete pair-production.

\section{Point Discontinuities and the Topo-Kinematic Projection Law}
\label{sec:projection_law}

In classical quantum mechanics, particle scattering and structural crossing interactions are governed strictly by probabilistic scattering matrices. Under the AVE Topo-Kinematic framework, intersection events across the discrete hardware of the vacuum graph are strictly resolved as \textbf{macroscopic signal routing operators}. When a wave traverses across nested configurations on the $K_4$ discrete graph, local phase interactions adhere exclusively to purely mechanical projection dynamics.

Beyond bulk impedance regime boundaries ($S(r)$ saturation), the discrete LC array mandates strict orientational coupling limits at geometric junctions. Because the discrete $0_1$ topological manifold acts as a phase-locked physical string supporting longitudinal stress parameters, routing transverse pressure waves identically mirrors macroscopic transmission line bending operations. 

Whenever a guided tracking wave alters its dimensional trajectory or intersects across a structured geometry intersecting at a macroscopic angle $\theta$, only the projection of the incident continuous wave mapped precisely onto the adjacent tracking axis survives the dynamic coupling:

\begin{resultbox}{The Topo-Kinematic Projection Law}
\begin{equation}
    Y_{loss} = c \frac{1 - \cos\theta}{2\pi^2}
\end{equation}
\end{resultbox}

This single universal continuous formula accurately defines spatial mechanical energy radiation boundaries across every physical scale defined in the AVE hierarchy organically:
\begin{itemize}
    \item \textbf{Subatomic Orbitals (Volume II):} Governs the geometric Phase C subshell scattering bounds where strings intersect across overlapping torus limits ($c_{intersections} \times \frac{1 - \cos\theta}{2\pi^2}$). When an outer tracking array encounters an fully-packed orthogonal topological continuous string ($l_{inner} > l_{outer}$, e.g. the $d^{10}$ shell), it encounters a \textbf{Polar Conjugate Mirror}. The rigid density of the inner boundary explicitly creates a Total Internal Reflection (TIR) boundary ($\Gamma < 0$) causing the outer tracking path to perfectly bypass intersecting drag ($Y_{loss} = 0.0$), structurally forcing dynamic volumetric Energy Compression over the valence array limit ($E_{val} = E_{base} \times (\text{knots} + 1)$).
    \item \textbf{Biological Proteins (Volume V):} Defines the radiation limits extracting capacitive binding fractions out of tight Amide-V biological backbone turns ($C_{bend} = \frac{1 - \cos\theta_i}{2\pi^2}$).
\end{itemize}

At a topological junction evaluating angle $\theta$, the uncoupled component natively extracted dynamically ($1 - \cos\theta$) radiates entirely into the surrounding topological grid as an environmental admittance to ground ($Y_{bend}$), explicitly eliminating local phase coherence and replacing classical statistical "probability loss" with explicit Ohmic/Joule mechanical transmission limits.

\section{Quantum Entanglement: The Phase-Locked Topological Thread}
\label{sec:entanglement}
A central result of standard Quantum Mechanics is ``Non-Locality''---the observation that two entangled particles exhibit correlated measurement outcomes regardless of spatial separation, as formalised by Bell's Theorem and confirmed by Clauser-Horne-Shimony-Holt (CHSH) inequality violations ($|S| = 2\sqrt{2} > 2$).

The AVE framework provides a \textbf{physical mechanism} for this correlation. Rather than treating non-locality as an abstract postulate, the lattice architecture reveals that entangled particles are connected by a \textit{topological thread}---a quantised phase winding on the $\mathcal{M}_A$ graph that functions as a phase-locked gear train, mechanically identical to the Meissner effect derived in Volume~III, Chapter~9.

\subsection{The Topological Thread}
When an entangled pair (e.g., $e^-e^+$) is synthesised from a single deformation of $\mathcal{M}_A$, two opposite-twist topological defects form. By conservation of topological charge on a connected graph (Axiom~1), the total phase winding along any path from particle~A to particle~B is quantised:
\begin{resultbox}{Quantised Phase Winding}
\begin{equation}
    \Delta\phi_{A \to B} = 2\pi \qquad (\text{winding number} = 1)
\end{equation}
\end{resultbox}

\noindent The thread connecting the two particles is a path of $N = d/\ell_{node}$ nodes on the $K_4$ graph, where $d$ is the spatial separation. Each node is an LC element (Axiom~1: no resistance $\Rightarrow$ lossless), and the two particle cores provide perfect short-circuit boundaries ($\Gamma = -1$, from Axiom~4 saturation). The thread therefore constitutes a \textbf{lossless short-short resonator}:
\begin{itemize}
    \item Characteristic impedance: $Z_0 = \sqrt{\mu_0/\epsilon_0} \approx 377 \; \Omega$
    \item Propagation speed: $v = c$
    \item Attenuation: $\alpha = 0$ (LC network, no dissipative element)
    \item Quality factor: $Q = \infty$
    \item Fundamental standing wave mode: $f_1 = c/(2d)$, \; $E_1 = \hbar\pi c / d$
\end{itemize}

\noindent The mode energy is \textbf{anti-confining}: $E_1 \propto 1/d$. Unlike the QCD flux tube ($E \propto d$, which confines quarks), the entanglement thread becomes \textit{energetically lighter} as the particles separate. For separations $d \gtrsim 1\;\mu\text{m}$, the mode energy drops below the thermal energy ($E_1 < k_B T$ at 300~K), and the thread's persistence transitions from energetic to \textbf{purely topological}---it survives not because it stores significant energy, but because destroying its winding requires pair creation (\S\ref{subsec:topo_protect}).

As the particles separate, the phase advance per node decreases:
\begin{equation}
    \delta\phi_{\text{per node}} = \frac{2\pi}{N} = \frac{2\pi \ell_{node}}{d} \xrightarrow{d \to \infty} 0
\end{equation}
The thread becomes \textbf{locally invisible} (each node is barely displaced from its ground state) but \textbf{globally present} (the cumulative winding is exactly $2\pi$).

\begin{figure}[h]
    \centering
    \includegraphics[width=1.0\textwidth]{entanglement_thread_energy.png}
    \caption{\textbf{Entanglement Thread Mode Properties (Simulation Output).} (Left) Mode energy $E_1 = \hbar\pi c/d$ vs.~separation. Three regimes are visible: energy-dominated ($d < 1\;\mu$m, red), comparable ($\sim 1\;\mu$m, yellow), and topology-dominated ($d > 1\;$mm, green). The dashed line marks $k_B T$ at 300~K. (Right) Phase advance per node $\delta\phi$ approaches zero as $d \to \infty$, while the total winding $\sum \delta\phi = 2\pi$ remains invariant.}
    \label{fig:thread_energy}
\end{figure}

\subsection{Topological Protection from Decoherence}
\label{subsec:topo_protect}
Thermal noise (the quantum foam of \S3.4) adds random phase perturbations $\delta_i$ at each node. However, because the fluctuations have zero mean, $\langle \sum_i \delta_i \rangle = 0$, the total winding remains $2\pi$. To destroy the thread, a coherent perturbation must shift the total winding by exactly $2\pi$---equivalent to spontaneous pair creation along the thread. The probability of this event is exponentially suppressed:
\begin{equation}
    P_{\text{break}} \sim \exp\!\left(-\frac{2m_e c^2}{k_B T}\right) \approx \exp(-3.95 \times 10^7) \approx 0 \qquad (T = 300 \text{ K})
\end{equation}
The thread is topologically immune to thermal decoherence at all physical temperatures.

\subsubsection{Falsifiable Prediction: Decoherence Regime Boundary}
The protection energy $2m_e c^2 \approx 1.022$~MeV defines a characteristic temperature:
\begin{equation}
    T_{\text{pair}} = \frac{2m_e c^2}{k_B} \approx 1.19 \times 10^{10} \text{ K}
\end{equation}
Below $T_{\text{pair}}$, the exponent $2m_e c^2 / k_B T \gg 1$ and the thread is absolutely stable. The transition from protected to vulnerable occurs sharply near $T \sim 10^9$--$10^{10}$~K, where the Boltzmann factor transitions from $\approx 0$ to $\approx 1$.

This constitutes a \textbf{falsifiable prediction} distinguishing AVE from standard quantum mechanics:
\begin{itemize}
    \item \textbf{AVE:} Entanglement decoherence has a sharp, temperature-dependent onset at the pair-creation threshold ($T \sim 10^9$~K). Below this temperature, decoherence proceeds only through environmental coupling to the thread endpoints.
    \item \textbf{Standard QM:} Decoherence is governed by environmental coupling strength alone, with no intrinsic temperature threshold tied to $2m_e c^2$.
\end{itemize}
This prediction is testable in quark-gluon plasma environments and heavy-ion collision experiments, where temperatures approach $10^{12}$~K.

\begin{figure}[h]
    \centering
    \includegraphics[width=0.7\textwidth]{entanglement_decoherence_regime.png}
    \caption{\textbf{Decoherence Regime Boundary (Simulation Output).} The thread-breaking probability $P_{\text{break}}$ vs.~temperature. Three regimes: stable (green, $T < 10^9$~K), transitional (yellow, $10^9$--$10^{10}$~K), and vulnerable (red, $T > 10^{10}$~K). The vertical line marks $T_{pair} = 2m_ec^2/k_B$.}
    \label{fig:decoherence_regime}
\end{figure}

\subsection{The Phase-Locked Gear Train}
The standing wave mode on the entanglement thread locks the two particle endpoints in anti-phase oscillation. This phase-lock is structurally identical to the Meissner mechanism derived in Volume~III, Chapter~9, where macroscopic superconductivity emerges from classical Kuramoto synchronisation of topological inductors.

\begin{center}
\small
\renewcommand{\arraystretch}{1.2}
\begin{tabular}{lll}
\toprule
\textbf{Property} & \textbf{Superconductor (Vol.~III)} & \textbf{Entanglement Thread} \\
\midrule
Oscillators       & $N$ electrons in wire       & $N$ LC nodes along thread \\
Phase lock         & Kuramoto below $T_c$        & Winding constraint ($\Delta\phi = 2\pi$) \\
Collective state   & Rigid gear train            & Rigid phase winding \\
Local perturbation & Applied B-field             & Ohmic measurement \\
Response           & Bulk inertia resists (Meissner) & Winding constrains partner \\
Operator           & $S(T/T_c) = \sqrt{1-(T/T_c)^2}$ & $S(\varepsilon/\alpha) = \sqrt{1-(\varepsilon/\alpha)^2}$ \\
\bottomrule
\end{tabular}
\end{center}

\noindent When Alice's detector couples to particle~A (an Ohmic load applying localised torque), she is not probing an isolated defect: she is turning one gear in a locked train. The winding constraint of the entire chain---a global topological invariant of the connected graph---constrains the outcome at particle~B.

\subsubsection{Universal $\Gamma \to -1$ Translation Matrix}
The entanglement thread is one instance of a universal mechanism: every phenomenon in the AVE framework where a boundary saturates to $\Gamma = -1$ is the same operator (Axiom~4) applied to a different sector. The complete cross-discipline isomorphism is:

\begin{center}
\scriptsize
\renewcommand{\arraystretch}{1.3}
\begin{tabularx}{\textwidth}{@{}l X X X X@{}}
\toprule
\textbf{Property} & \textbf{Particle Confinement} & \textbf{Plasma Cutoff} & \textbf{Superconductor} & \textbf{Entanglement Thread} \\
\midrule
Domain            & Topology (Vol.~I)  & EM (Vol.~I)            & Condensed matter (Vol.~III)  & Quantum info (Vol.~I) \\
Saturated field   & $\varepsilon(r)$   & $\varepsilon_{eff}$    & $\mu_{eff}$                  & $Z_{core}$ \\
Saturation input  & $r/\ell_{node}$    & $V/V_{snap}$           & $B/B_c$                      & $\varepsilon/\alpha$ \\
Operator          & $S = \sqrt{1-r^2}$ & $S = \sqrt{1-(V/V_s)^2}$ & $S = \sqrt{1-(B/B_c)^2}$ & $S = \sqrt{1-(\varepsilon/\alpha)^2}$ \\
$\Gamma \to -1$ means & Particle trapped  & EM wave reflected      & B-field expelled             & Winding topologically locked \\
Characteristic scale & $\ell_{node}$   & $\delta = c/\omega_p$  & $\lambda_L$                  & $d$ (separation) \\
Collective state  & Knot (stable defect) & Plasma cutoff         & Rigid gear train             & Phase-locked thread \\
Engine module     & \texttt{faddeev\_skyrme} & \texttt{plasma/cutoff} & \texttt{plasma/}\allowbreak\texttt{superconductor} & \texttt{topological/}\allowbreak\texttt{entanglement\_thread} \\
\bottomrule
\end{tabularx}
\end{center}

\noindent All four phenomena are evaluated by the same function \texttt{saturation\_factor()} from \texttt{ave.\allowbreak axioms.\allowbreak scale\_\allowbreak invariant}. There is no domain-specific physics: the saturation operator is universal, and the ``quantum non-locality'' of entanglement is structurally identical to the ``classical non-locality'' of the Meissner effect.

\subsection{Deriving the Angular Correlation Function}
\label{subsec:bell_correlation}
The angular dependence of the entanglement correlation derives from three AVE ingredients:

\begin{enumerate}
    \item \textbf{Spin-1/2 as M\"obius topology (Axiom~1).} The electron unknot ($0_1$) is a Beltrami standing wave ($\nabla \times \mathbf{A} = k\mathbf{A}$) on the chiral $K_4$ graph. The M\"obius-strip topology of the chiral labyrinth requires $720^\circ$ for a complete cycle, producing a physical \textbf{half-angle coupling}: when a detector axis $\hat{a}$ makes angle $\theta$ with the defect's rotation axis $\hat{n}$, the effective phase mismatch is $\theta/2$.

    \item \textbf{Binary outcome from saturation (Axiom~4).} The particle's saturated boundary ($\Gamma = -1$) supports a standing wave with exactly two antinodes. The detector couples to whichever antinode has lower impedance mismatch, producing a binary ($\pm$) outcome.

    \item \textbf{Probability from Ohmic extraction (Axiom~3).} The Born rule, derived in \S3.5, gives the probability of each outcome from the Joule power extracted by the detector:
    \begin{equation}
        P(+\hat{a}) = \cos^2\!\left(\frac{\theta}{2}\right), \qquad P(-\hat{a}) = \sin^2\!\left(\frac{\theta}{2}\right)
    \end{equation}
\end{enumerate}

\noindent After Alice measures $+\hat{a}$, the gear-train constraint sets particle~B to state $-\hat{a}$ (anti-correlation from winding conservation). Bob then measures along $\hat{b}$, at angle $(\pi - \theta_{ab})$ from $-\hat{a}$:
\begin{align}
    P(B{=}+ \mid A{=}+) &= \sin^2(\theta_{ab}/2) \\
    P(B{=}- \mid A{=}+) &= \cos^2(\theta_{ab}/2)
\end{align}
The correlation function evaluates to:
\begin{resultbox}{Bell Angular Correlation}
\begin{equation}
    E(\hat{a}, \hat{b}) = -\cos\theta_{ab}
\end{equation}
\end{resultbox}

\subsection{CHSH Inequality Violation}
The Clauser-Horne-Shimony-Holt parameter is defined as $S = E(a,b) - E(a,b') + E(a',b) + E(a',b')$. With $E = -\cos\theta$ and optimal detector angles separated by $\delta$:
\begin{equation}
    S(\delta) = -3\cos\delta + \cos 3\delta
\end{equation}
Maximised at $\delta = \pi/4$:
\begin{resultbox}{CHSH Violation}
\begin{equation}
    |S|_{\max} = 2\sqrt{2} \approx 2.828
\end{equation}
\end{resultbox}
\noindent This exceeds the classical bound of $2$ and matches the Tsirelson bound exactly, confirming full quantum-mechanical angular correlations from AVE first principles.

\begin{figure}[h]
    \centering
    \includegraphics[width=1.0\textwidth]{entanglement_bell_chsh.png}
    \caption{\textbf{Bell Correlation and CHSH Violation (Simulation Output).} (Left) Angular correlation function $E(\theta) = -\cos\theta$, showing anti-correlation (blue) transitioning to correlation (red). (Right) CHSH parameter $|S(\delta)|$ vs.~detector spacing. The cyan region exceeds the classical bound ($|S| = 2$, red dashed), reaching the Tsirelson bound $2\sqrt{2}$ (gold dotted) at $\delta^* = 45°$. Both curves are computed entirely from Axioms~1--4 with no imported quantum postulates.}
    \label{fig:bell_chsh}
\end{figure}

\begin{examplebox}[Entanglement Thread at Key Separations]
\textbf{Problem:} Compute the thread mode energy, frequency, and phase advance per node for an entangled pair at nuclear, laboratory, and astronomical separations.

\textbf{Solution:} Using the engine functions $E_1 = \hbar\pi c/d$, $f_1 = c/(2d)$, and $\delta\phi = 2\pi\ell_{node}/d$:

\begin{center}
\small
\renewcommand{\arraystretch}{1.2}
\begin{tabular}{lrrrr}
\toprule
\textbf{Separation} & \textbf{$f_1$} & \textbf{$E_1$} & \textbf{$\delta\phi$/node} & \textbf{Regime} \\
\midrule
$\ell_{node}$ (electron) & $3.88 \times 10^{20}$ Hz & 1.61 MeV & $360°$ & Energy \\
1~pm (atomic) & $1.50 \times 10^{20}$ Hz & 620 keV & $1.4 \times 10^{-1\circ}$ & Energy \\
1~$\mu$m & $1.50 \times 10^{14}$ Hz & 0.62 eV & $1.4 \times 10^{-4\circ}$ & Comparable \\
1~m & $1.50 \times 10^{8}$ Hz & $6.2 \times 10^{-7}$ eV & $1.4 \times 10^{-10\circ}$ & Topological \\
1~km & $1.50 \times 10^{5}$ Hz & $6.2 \times 10^{-10}$ eV & $1.4 \times 10^{-13\circ}$ & Topological \\
\bottomrule
\end{tabular}
\end{center}

At laboratory separations ($d \gtrsim 1$~m), the mode energy is $\sim 10^5$ times smaller than the thermal energy. The thread persists because destroying the $2\pi$ winding requires $2m_ec^2 \approx 1$~MeV---a topological barrier, not an energetic one.
\end{examplebox}

\subsection{No-Signalling Theorem}
\label{subsec:no_signaling}
Bob's marginal probability is independent of Alice's detector setting $\hat{a}$:
\begin{equation}
    P(B{=}+) = \tfrac{1}{2}\sin^2(\theta_{ab}/2) + \tfrac{1}{2}\cos^2(\theta_{ab}/2) = \tfrac{1}{2}
\end{equation}
Bob always observes a 50/50 outcome distribution. Alice cannot control which outcome she receives (the gear-train response is stochastic from her perspective due to the thermal noise floor), so no encodable information traverses the thread. The no-signalling theorem holds.

\subsection{K4-TLM Lattice Verification}
\label{subsec:k4tlm_entanglement}
To verify the topological protection mechanism at the lattice level, a full 3D K4-TLM simulation was run on a $32^3$-node Diamond lattice with nonlinear Axiom~4 saturation enabled. Two saturated defect walls ($\Gamma = -1$, modelling particle cores) were placed 19 nodes apart, connected by a $2\pi$ phase-wound standing wave initialised with a $\sin(\pi x/L)$ envelope and helical port structure.

Three scenarios were compared:
\begin{enumerate}
    \item \textbf{Vacuum (T = 0):} No perturbations. The thread's phase coherence evolves from 1.00 to 0.73 as the evanescent lateral tail spreads on the periodic lattice---a finite-size boundary effect, not decoherence.
    \item \textbf{Mild noise ($T \ll T_{pair}$):} Random Gaussian perturbations ($\sigma = 0.005$) applied to all active nodes. Coherence: $1.00 \to 0.74$---\textbf{indistinguishable from vacuum}. The winding is topologically immune to sub-threshold noise.
    \item \textbf{Extreme noise ($T \gg T_{pair}$):} Continuous high-amplitude perturbations ($\sigma = 0.3$) overwhelm the saturation barrier. Coherence: $1.00 \to 0.28$---the $2\pi$ winding is destroyed.
\end{enumerate}

\noindent The key result is that Scenarios~1 and~2 are statistically identical: the $2\pi$ winding is a topological invariant of the lattice, and perturbations below the pair-creation energy cannot unwrap it. Only when the noise amplitude exceeds the Axiom~4 saturation threshold does the thread decohere.

\begin{figure}[h]
    \centering
    \includegraphics[width=1.0\textwidth]{k4_tlm_decoherence_3d.png}
    \caption{\textbf{K4-TLM 3D Decoherence (Simulation Output).} Final-state energy density on the Diamond lattice for the three scenarios. Yellow stars mark the saturated defect cores ($\Gamma = -1$). Cyan nodes carry the phase-wound thread; magenta indicates scattered noise energy. In the vacuum and mild-noise cases, the thread structure is preserved. In the extreme-noise case, the winding is disrupted and energy fills the lattice isotropically.}
    \label{fig:k4tlm_decoherence}
\end{figure}

\begin{figure}[h]
    \centering
    \includegraphics[width=0.8\textwidth]{k4_tlm_winding_evolution.png}
    \caption{\textbf{Thread Winding Evolution (Simulation Output).} Phase coherence (top) and normalised amplitude (bottom) vs.~timestep. The vacuum (cyan) and mild-noise (orange) traces overlap, confirming topological protection. The extreme-noise trace (red) decays sharply as the winding is overwhelmed.}
    \label{fig:k4tlm_winding}
\end{figure}

\section*{Chapter Summary}
\begin{itemize}
    \item The Lagrangian of QED is mechanically isomorphic to the continuous inductive bulk stress of the $\mathcal{M}_A$ lattice, where the Vector Potential ($\mathbf{A}$) maps directly to mass flow rate.
    \item Standard quantum artifacts like complex Hilbert spaces and the Uncertainty Principle are natural signal-processing consequences of sending finite-bandwidth data across a discrete, Nyquist-limited sampling grid.
    \item ``Quantum Foam'' is simply the irreducible baseline RMS thermal noise floor of the universe's hardware substrate. Virtual particles are failed topological knots that dissipate before achieving resonant stability.
    \item Quantum Entanglement arises from a topological thread---a quantised $2\pi$ phase winding on the $K_4$ graph forming a lossless short-short resonator. The thread acts as a phase-locked gear train, structurally identical to the Meissner mechanism of superconductivity (Volume~III, Chapter~9). The M\"obius half-angle and Axiom~4 binary quantiser produce the angular correlation $E(\hat{a},\hat{b}) = -\cos\theta_{ab}$, yielding CHSH $= 2\sqrt{2}$.
\end{itemize}

\begin{summarybox}
\begin{itemize}
    \item The Schr\"odinger Equation and Generalised Uncertainty Principle emerge from discrete-bandwidth LC circuit resonance on the $\mathcal{M}_A$ lattice.
    \item Wave--particle duality is the interference pattern of finite-bandwidth lattice signals; quantum foam is baseline electrical noise.
    \item Quantum entanglement is mediated by a topological thread---a phase-locked gear train in the vacuum lattice. The correlation $E = -\cos\theta_{ab}$ and the CHSH violation $|S| = 2\sqrt{2}$ are derived from M\"obius half-angle geometry (Axiom~1), dielectric saturation quantisation (Axiom~4), and the Ohmic Born rule (Axiom~3), with no imported quantum postulates.
\end{itemize}
\end{summarybox}

\begin{exercisebox}
\begin{enumerate}
    \item Using the dimensional definition of the Vector Potential as Mass Flow ($\mathbf{A} \equiv \xi_{topo}^{-1}\, [\text{kg/s}]$), show that the magnetic field $\mathbf{B} = \nabla \times \mathbf{A}$ corresponds dimensionally to macroscopic fluid vorticity.
    \item Verify the CHSH bound: with $E(\hat{a},\hat{b}) = -\cos\theta_{ab}$ and detector angles $a=0$, $b=\pi/4$, $a'=\pi/2$, $b'=3\pi/4$, compute $S$ and confirm $|S| = 2\sqrt{2}$.
    \item The entanglement thread has mode energy $E_1 = \hbar\pi c/d$. Compute the mode energy for a pair separated by $d = 1$~km and compare to the thermal energy $k_BT$ at $T = 300$~K. Why does the thread persist despite $E_1 \ll k_BT$?
\end{enumerate}
\end{exercisebox}